% Ethics / responsible-release section (task 10.2)

\section{Ethics and Responsible Release}
\label{sec:ethics}

\subsection{Dual-Use Considerations}
A detection system for RAG poisoning is inherently dual-use: understanding
how detectors fail is also, necessarily, a partial recipe for evading them.
Our own adversarial red-teaming (Section~\ref{sec:results}) surfaced a real
evasion technique -- framing a malicious instruction as generic
documentation/license boilerplate to dilute the embedding-anomaly signal
(``boilerplate anchoring''). We report this finding because withholding it
would leave practitioners unaware of a real blind spot in a published
class of defenses, which is a worse outcome than disclosing it alongside
a discussion of mitigations (Section~\ref{sec:limitations}). This follows
the standard responsible-disclosure norm in security research: publish the
weakness together with what is known about addressing it, rather than
publishing neither or publishing the weakness alone.

\subsection{Data Provenance and Consent}
\begin{itemize}
  \item Clean baseline documents (\texttt{data/clean/}) are drawn from
    real, public open-source project READMEs; sources and licenses are
    logged in \texttt{data/\_provenance/}.
  \item Poisoned and adversarial documents (\texttt{data/poisoned/},
    \texttt{data/adversarial/}) were authored by the team specifically for
    this evaluation. They are not drawn from, nor do they reproduce, any
    real production attack, exploit, or incident.
  \item No real user data, private corpora, or personally identifying
    information is used anywhere in the dataset.
\end{itemize}

\subsection{Scope of Claims}
Consistent with \texttt{related\_work.md}, this project does not claim to
introduce a novel detection algorithm. Its contribution is a comparative
failure, blind-spot, and latency analysis of three published detection
signals (instruction-pattern matching, embedding anomaly, perplexity
scoring) and their fusion. We are explicit about this scope so the work is
not mistaken for -- or oversold as -- a new, patent-worthy detection
technique in isolation from calibration method (Phase 5's conformal
calibration work is the project's primary candidate for a genuinely novel,
disclosable technical contribution; see \texttt{docs/prior\_art.md}).

\subsection{Pre-Publication Review}
Per the project roadmap (task 11.1), this work will be reviewed by the
institution's IP/patent cell \emph{before} any public posting of code,
data, or the paper, to confirm what can safely be disclosed and when. No
component of this project has been publicly posted outside the private
team repository as of this writing.

\subsection{Intended Use}
PDS is intended to help practitioners evaluate and compare RAG-poisoning
defenses, not to serve as a how-to guide for evading them. Where we discuss
a successful evasion, we pair it with the specific gap it exposes
(Section~\ref{sec:limitations}) and the concrete next step planned to close
it (e.g., the credential-exposure pattern added in the dataset in response
to the \texttt{adv\_010} finding).
